% !TeX spellcheck = ru_RU
% !TEX root = vkr.tex

\section{Описание решения}
\label{sec:solution}

\subsection{Требования к системе}

Требования сформулированы совместно с консультантом (м.н.с.\ лаборатории Гербарий ГБС РАН, Шкурко~А.~В.)
на основе технического задания.

\subsubsection*{Требования заказчика}

\begin{enumerate}
    \item Автоматическое распознавание формы листа и построение полигональной маски контура.
    \item Определение длины и ширины полигона.
    \item Расстановка точек вдоль перпендикуляров к главной оси с заданным шагом (например, $\frac{1}{2}$ или $\frac{1}{4}$ длины листа).
    \item Просмотр исходного изображения и маски рядом на экране.
    \item Возможность задать масштабный коэффициент ($1~\text{пиксель} = \ldots~\text{мм}$) для вывода измерений в физических единицах.
    \item Экспорт результатов в табличный формат (CSV / Excel).
    \item Пакетная обработка серий изображений.
    \item Возможность ручной корректировки автоматически построенной маски.
\end{enumerate}

\subsubsection*{Функциональные требования}

На основе технического задания выделены следующие варианты использования:

\begin{itemize}
    \item \textbf{UC1.} Загрузить набор изображений для обработки.
    \item \textbf{UC2.} Получить маску сегментации по интерактивным подсказкам (точки, ограничивающий прямоугольник).
    \item \textbf{UC3.} Скорректировать маску вручную в режиме кисти/ластика.
    \item \textbf{UC4.} Добавить подтверждённую маску в буфер или удалить её.
    \item \textbf{UC5.} Задать масштабный коэффициент пиксель~→~мм.
    \item \textbf{UC6.} Запустить пакетное измерение всех сохранённых масок.
    \item \textbf{UC7.} Экспортировать таблицу результатов в~CSV.
\end{itemize}

\subsubsection*{Нефункциональные требования}

\begin{itemize}
    \item Работа на CPU без выделенного GPU~--- приложение должно запускаться на стандартных рабочих станциях лаборатории.
    \item Время отклика модели сегментации не более 3~с на одно изображение.
    \item Поддержка форматов JPEG, PNG, TIFF.
    \item Воспроизводимость результатов: одни и те же входные данные и подсказки дают одинаковый вывод.
    \item Единый формат хранения масок (PNG) и измерений (CSV) для последующей статистической обработки.
\end{itemize}

\subsection{Архитектура приложения}

Приложение реализовано на языке Python с использованием фреймворка PyQt6 для графического интерфейса
и библиотеки PyTorch для инференса модели MobileSAM.
Исходный код опубликован по адресу \url{https://github.com/pavelrubanov/segmentation_system}.

Система состоит из двух слоёв: \textit{ядра} (core) и \textit{пользовательского интерфейса} (ui).
Ядро не зависит от GUI и включает обёртку над моделью сегментации, модель данных изображений
и модуль измерений. Слой UI реализует окна и виджеты PyQt6.

\begin{figure}[H]
    \centering
    % TODO: вставить System architecture.png
    % \includegraphics[width=0.9\linewidth]{figures/architecture.png}
    \caption{Архитектура системы}
    \label{fig:architecture}
\end{figure}

Точка входа~--- \texttt{main.py}~--- принимает аргумент \texttt{-{}-checkpoint} с путём к весам модели,
проверяет его наличие и запускает главное окно навигации (\texttt{NavigationWindow}).

\subsection{Диаграмма активностей}

На рисунке~\ref{fig:activity} представлен типовой рабочий цикл биолога при работе с системой.
Ключевой момент~--- решение об отборе листа (ромб «Лист пригоден?») остаётся за экспертом,
тогда как система берёт на себя сегментацию и все последующие вычисления.

\begin{figure}[H]
\centering
\begin{tikzpicture}[
    font=\footnotesize,
    node distance=0.55cm,
    action/.style={rectangle, draw, rounded corners=3pt, align=center,
                   minimum width=4.8cm, minimum height=0.75cm, fill=white},
    decision/.style={diamond, draw, align=center, inner sep=1pt,
                     minimum width=1.0cm, fill=white},
    arr/.style={->, >=Stealth, thick},
    label/.style={font=\scriptsize, inner sep=1pt},
]

% --- nodes ---
\node[draw, circle, fill=black, minimum size=0.28cm] (start) {};

\node[action,  below=of start]       (load)       {Загрузить папку с изображениями};
\node[action,  below=of load]        (nav)        {Перейти к следующему изображению};
\node[action,  below=of nav]         (click)      {Кликнуть на лист\\(positive point / bbox)};
\node[action,  below=of click]       (sam)        {\textit{[система]} Построить маску (MobileSAM)};
\node[decision,below=of sam]         (ok)         {Маска\\точная?};
\node[action,  below=of ok]          (save)       {\textit{[система]} Сохранить маску в буфер};
\node[decision,below=of save]        (moreleaves) {Ещё листья\\на изображении?};
\node[decision,below=of moreleaves]  (moreimages) {Ещё\\изображения?};
\node[action,  below=of moreimages]  (batch)      {Запустить пакетные измерения};
\node[action,  below=of batch]       (export)     {Экспортировать CSV};

\node[draw, circle, fill=black, minimum size=0.28cm,
      below=of export] (end) {};
\node[draw, circle, minimum size=0.38cm] at (end) {};  % двойной круг

% боковой узел: правка маски
\node[action, right=1.8cm of ok] (edit) {Исправить маску\\(кисть / ластик)};

% --- arrows ---
\draw[arr] (start) -- (load);
\draw[arr] (load)  -- (nav);
\draw[arr] (nav)   -- (click);
\draw[arr] (click) -- (sam);
\draw[arr] (sam)   -- (ok);

% ok? → да ↓, нет → edit → сразу к save
\draw[arr] (ok)       -- node[label, left]{да} (save);
\draw[arr] (ok.east)  -- node[label, above]{нет} (edit.west);
\draw[arr] (edit.south) |- (save.east);

\draw[arr] (save) -- (moreleaves);

% moreleaves? → нет ↓, да → влево и вверх к click.west (внутренний контур)
\draw[arr] (moreleaves) -- node[label, left]{нет} (moreimages);
\draw[arr] (moreleaves.west) -- ++(-2.0,0) node[label, above left]{да}
           |- (click.west);

% moreimages? → нет ↓, да → влево и вверх к nav.west (внешний контур)
\draw[arr] (moreimages) -- node[label, left]{нет} (batch);
\draw[arr] (moreimages.west) -- ++(-3.0,0) node[label, above left]{да}
           |- (nav.west);

\draw[arr] (batch)  -- (export);
\draw[arr] (export) -- (end);

\end{tikzpicture}
\caption{Диаграмма активностей: рабочий цикл пользователя}
\label{fig:activity}
\end{figure}

\subsection{Диаграмма классов}

Ключевые классы системы:

\begin{itemize}
    \item \texttt{Predictor}~--- обёртка над MobileSAM.
          Метод \texttt{predict(points, labels, box)} принимает список координат точек,
          их меток ($+1$ / $-1$) и опциональный ограничивающий прямоугольник;
          возвращает бинарную маску в виде массива \texttt{uint8} (значения 0 и 255).
          Инференс выполняется на CPU с помощью \texttt{torch.no\_grad()}.

    \item \texttt{ImageWithMasks}~--- датакласс, хранящий изображение (\texttt{numpy} RGB) и имя файла.
          Маски не хранятся в памяти~--- только пути к файлам на диске, что снижает потребление RAM
          при пакетной обработке.

    \item \texttt{LeafMetrics}~--- результат измерения одного листа:
          \texttt{length\_px}, \texttt{width\_px} и список \texttt{SectionWidth}
          для семи поперечных сечений.

    \item \texttt{SegmentationCanvas}~--- виджет \texttt{QGraphicsView} с двумя режимами:
          \textit{PROMPTS} (клики мышью задают подсказки для SAM) и
          \textit{EDIT} (кисть/ластик для ручной коррекции маски).
          Поддерживает зум в точку курсора и панорамирование по \texttt{Ctrl+ЛКМ}.

    \item \texttt{MasksBuffer}~--- галерея миниатюр сохранённых масок (256$\times$256 px).
          Маски индексируются по имени изображения (\texttt{image\_mask\_0001.png} и т.д.).
\end{itemize}

\subsection{Диаграмма последовательности}

Типовой сценарий работы пользователя:

\begin{enumerate}
    \item Загрузить набор изображений~→ \texttt{NavigationWindow} открывает \texttt{SegmentationWindow}.
    \item Щёлкнуть по листу левой кнопкой мыши (positive point)~→
          \texttt{SegmentationCanvas} вызывает \texttt{Predictor.predict()}~→
          полупрозрачный overlay маски отображается поверх изображения.
    \item При необходимости переключиться в режим EDIT и подправить маску кистью.
    \item Подтвердить маску~→ файл PNG сохраняется на диск, \texttt{MasksBuffer} обновляет миниатюру.
    \item Перейти в \texttt{ProcessingWindow}, выбрать папку с масками, запустить измерения.
    \item \texttt{leaf\_measure.measure\_leaf()} обрабатывает каждую маску~→ CSV экспортируется.
\end{enumerate}

\subsection{Алгоритм параметрических измерений}

Модуль \texttt{leaf\_measure.py} реализует измерение геометрических характеристик листа
на основе анализа главных компонент (PCA) контура маски.

\begin{enumerate}
    \item Из бинарной маски извлекается внешний контур (OpenCV \texttt{findContours}).
    \item К координатам точек контура применяется PCA:
          первый главный вектор соответствует главной оси листа (длина),
          второй~--- перпендикулярному направлению (ширина).
    \item Длина листа вычисляется как проекция контура на первый главный вектор.
    \item Ширина в точке $t \in \{0.125, 0.25, 0.375, 0.5, 0.625, 0.75, 0.875\}$~---
          расстояние между двумя пересечениями перпендикуляра к главной оси с контуром.
    \item При наличии масштабного коэффициента значения пересчитываются в миллиметры.
\end{enumerate}

Для каждой обработанной маски генерируется изображение-визуализация (\texttt{*\_vis.png})
с нанесёнными контуром, осями и поперечными сечениями.

\subsection{Стек технологий}

\begin{table}[H]
\centering
\caption{Использованные технологии}
\label{tab:stack}
\begin{tabular}{lll}
\hline
\textbf{Компонент} & \textbf{Технология} & \textbf{Обоснование} \\
\hline
Модель сегментации & MobileSAM (PyTorch) & Foundation model, CPU, 39~МБ \\
Графический интерфейс & PyQt6 & Кроссплатформенный, мощный QGraphicsView \\
Обработка изображений & OpenCV, NumPy & Стандарт для CV-задач \\
Загрузка изображений & Pillow & TIFF, PNG, JPEG \\
Экспорт результатов & csv (stdlib) & Совместимость с Excel \\
\hline
\end{tabular}
\end{table}

\subsection{Тестирование}

Для верификации модуля измерений написаны модульные тесты:

\begin{itemize}
    \item \textit{Прямоугольная маска}: заданной ширины и высоты~---
          ожидаемые длина и ширина совпадают с расчётными с погрешностью менее 1\%.
    \item \textit{Эллиптическая маска}: проверяется симметричность профиля
          (ширины в точках $t$ и $1-t$ должны совпадать).
    \item \textit{Интеграционный тест}: загрузка реальных тестовых изображений из директории \texttt{test/images/},
          получение маски через \texttt{Predictor}, запуск \texttt{measure\_leaf()} и
          проверка формата вывода CSV.
\end{itemize}
